% Proposed focal addition; does not replace the existing electronic algebra.
% Insertion: after the proposition "Algebra of the local electronic block"
% and its explanatory paragraph, before "Transfer between sheets...".
% Provenance and scope: LOCALIZADORES_Y_ALCANCE.md, in this folder.
\subsection{Spinor representation of the electronic block and helicity}
\label{sec:el-representacion-espinorial}

The electronic structure has a representation preserving more
information than its scalar evaluation. The principal plane of the two
three-dimensional APP projectors, selected by the tritic mode
\((1,-1,0)^3\), has already provided the operators \(S\) and \(J\).
The former is a symmetric involution; the latter is the normalized
commutator and preserves sheet orientation. Their composition makes it possible
to construct explicitly a rotation representation and separate its
two helicity components. This construction acts on the electronic
block; the TPK route record and its memory remain as data
of the block's provenance.

Let \(E_e\) be the real principal plane, equipped with the inner product inherited
from the conditional expectations, and let
\(\mathscr S_e=E_e\otimes_{\mathbb R}\mathbb C\) be its complexification.
The operators already obtained satisfy
\[
 S^*=S,\qquad J^*=-J,\qquad
 S^2=I,\qquad J^2=-I,\qquad SJ=-JS.
\]
Complexification extends the representation field. It preserves the
real operators and allows the scalar unit \(i\) to be distinguished from the
endomorphism \(J\), although both square to \(-1\) in their
respective domains.

\begin{proposition}[Hermitian triple induced by the two sheets]
\label{prop:el-terna-hermitica}
The operators
\begin{equation}
 \sigma_1=SJ,\qquad \sigma_2=-iJ,\qquad \sigma_3=S
 \label{eq:el-terna-hermitica}
\end{equation}
are Hermitian and traceless. They satisfy
\begin{equation}
 \sigma_a\sigma_b
 =\delta_{ab}I+i\sum_{c=1}^3\epsilon_{abc}\sigma_c,
 \qquad
 \frac12\operatorname{tr}(\sigma_a\sigma_b)=\delta_{ab},
 \label{eq:el-algebra-terna}
\end{equation}
where \(\epsilon_{123}=1\). They therefore form an orthonormal basis
of the real space \(V_e\) of traceless Hermitian endomorphisms of
\(\mathscr S_e\), for the inner product
\(\langle X,Y\rangle=\frac12\operatorname{tr}(XY)\).
\end{proposition}
\begin{proof}
The preceding relations give \((SJ)^*=(-J)S=SJ\) and
\((-iJ)^*=-iJ\). The squares of all three matrices are \(I\).
Furthermore,
\[
 (SJ)(-iJ)=iS,\qquad (-iJ)S=iSJ,\qquad S(SJ)=J=i(-iJ).
\]
Reversing the order changes the sign. This proves the first identity
in \eqref{eq:el-algebra-terna}. In the orthonormal basis of the principal plane,
the matrices are
\[
 \sigma_1=\begin{pmatrix}0&1\\1&0\end{pmatrix},\quad
 \sigma_2=\begin{pmatrix}0&-i\\i&0\end{pmatrix},\quad
 \sigma_3=\begin{pmatrix}1&0\\0&-1\end{pmatrix}.
\]
Their trace is zero, and taking traces in the product identity proves
orthonormality. Every traceless Hermitian endomorphism has the form
\(\bigl(\begin{smallmatrix}z&x-iy\\x+iy&-z\end{smallmatrix}\bigr)\),
with real \(x,y,z\), and thus belongs to their real span.
\end{proof}

In the usual matrix representation, this triple is called the
Pauli matrices. Their use here follows the construction of
\(P,Q,S,J\): they are obtained through the products
\eqref{eq:el-terna-hermitica}. The direction space of this
representation is the unit sphere of \(V_e\), whose positive form
has just been constructed through the trace.

\begin{proposition}[Rotation generators and directional decomposition]
\label{prop:el-helicidad}
Define \(L_a=\sigma_a/2\). Then
\begin{equation}
 [L_a,L_b]=i\sum_c\epsilon_{abc}L_c,
 \qquad \sum_{a=1}^3L_a^2=\frac34 I.
 \label{eq:el-generadores-spin}
\end{equation}
For \(n=(n_1,n_2,n_3)\) with \(\sum_an_a^2=1\), let
\[
 H_e(n)=\sum_an_a\sigma_a,\qquad
 h_e(n)=\frac12H_e(n),\qquad
 \Pi_\pm(n)=\frac12\bigl(I\pm H_e(n)\bigr).
\]
The linear formula \(H_e(v)=\sum_av_a\sigma_a\) is also used
for \(v\) of arbitrary norm.
One has
\begin{equation}
 \mathscr S_e=\operatorname{im}\Pi_+(n)
             \mathbin{\oplus}\operatorname{im}\Pi_-(n),
 \qquad
 h_e(n)\Pi_\pm(n)=\pm\frac12\Pi_\pm(n).
 \label{eq:el-descomposicion-helicidad}
\end{equation}
The two summands are orthogonal complex lines. Rotation
transport around \(n\) is
\begin{equation}
 U_n(\theta)=\exp\!\left(-\frac{i\theta}{2}H_e(n)\right)
 =\cos\frac\theta2\,I-i\sin\frac\theta2\,H_e(n),
 \qquad
 U_n(2\pi)=-I,\quad U_n(4\pi)=I.
 \label{eq:el-retorno-spin}
\end{equation}
\end{proposition}
\begin{proof}
The product identity of the triple gives the commutators in
\eqref{eq:el-generadores-spin}; their squares sum to \(3I/4\).
Antisymmetric terms cancel when multiplying
\(H_e(n)\) by itself, so that \(H_e(n)^2=I\).
It follows that
\[
 \Pi_\pm(n)^2=\Pi_\pm(n),\qquad
 \Pi_+(n)\Pi_-(n)=0,\qquad \Pi_+(n)+\Pi_-(n)=I.
\]
They are Hermitian and have trace \(1\); their rank is therefore \(1\).
Multiplication by \(H_e(n)/2\) proves
\eqref{eq:el-descomposicion-helicidad}. The exponential series
splits into even and odd powers because \(H_e(n)^2=I\), producing
\eqref{eq:el-retorno-spin}. It also directly proves
unitarity and the law \(U_n(\theta+\psi)=U_n(\theta)U_n(\psi)\).

Angular normalization is verified on \(V_e\). For every
\(v\in\mathbb R^3\), the same product identity gives
\[
 U_n(\theta)H_e(v)U_n(\theta)^*
 =H_e\!\left(v\cos\theta+(n\mathbin{\times}v)\sin\theta
       +n(n\mathbin{\cdot}v)(1-\cos\theta)\right).
\]
The induced action is a rotation of angle \(\theta\) and axis \(n\).
The factor \(1/2\) in the spinor generator is therefore the one
corresponding to that angular parameterization. A complete rotation restores
the direction in \(V_e\) and changes the sign of the vector in
\(\mathscr S_e\); two complete rotations restore both.
\end{proof}

The representation is irreducible: a complex subspace invariant under
the three matrices would be invariant under all of \(M_2(\mathbb C)\), the
algebra they generate. Its only invariant subspaces are \(0\) and
\(\mathscr S_e\). The relations
\eqref{eq:el-generadores-spin} thus identify the spin-\(1/2\)
representation. If \(n\) is the propagation direction in a
kinematic realization, \(h_e(n)\) is its dimensionless helicity, with
values \(\pm1/2\). The kinematic direction is declared when that
realization is made; the two projectors are defined for every direction
before a particular propagation state is selected.

\paragraph{Preservation of the electronic structure and types of inversion.}
Direction inversion has an exact law:
\[
 H_e(-n)=-H_e(n),\qquad \Pi_\pm(-n)=\Pi_\mp(n).
\]
This operation exchanges helicity labels relative to a
direction. Charge conjugation, by contrast, acts on the oriented
signature of the material route; cellular inversion acts on the
atlas coordinates. Their domains and actions remain distinct.
Uniqueness of cell \((5,5)\) is compatible with the two lines in
\eqref{eq:el-descomposicion-helicidad}: an atlas cell is not a
vector in the spinor representation it carries.

The electronic evaluation \(\varepsilon_e^\beta\) preserves its
formula and readers. On \(\mathscr S_e\), its scalar realization
is \(\varepsilon_e^\beta I\), which commutes with the generators and with
\(\Pi_\pm(n)\). Helicity decomposition adds rotational
information without modifying the evaluation or turning spin into a
multiplicative mass coefficient.

Chirality corresponds to the grading of the relativistic representation
used: an involution \(\Gamma_{\rm ch}\) determines its
projectors \((I\mp\Gamma_{\rm ch})/2\). The helicity in
\eqref{eq:el-descomposicion-helicidad} additionally depends on \(n\).
For its part, the holonomy of a real Möbius band describes the change
of sign of a line when traversing a loop in its base. The law of that line,
the spinor return \eqref{eq:el-retorno-spin}, and nonadic phase holonomy
with memory advance are transport laws on specified
objects. Their comparison preserves those objects: the helical form
of a trajectory does not replace the representation just
constructed or by itself fix chirality or charge.
